\section{Numerical classes, psh functions and the envelope property on a projective variety}
\label{sec:source-numerical-psh}

Throughout, $k$ is a field with the trivial absolute value and $p:X\to\operatorname{Spec}k$ is
projective; from Definition~\ref{def:var-psh} on, $X$ is integral. We write $N^1(X)_\R$ for the
real Picard space of $X$ modulo the classes of degree zero on every closed integral curve
(Theorem~\ref{thm:native-real-numerical-classes}), $[L]$ for the class of an invertible sheaf
$L$, $\operatorname{Amp}(X)\subset N^1(X)_\R$ for the convex cone generated by the classes of
ample invertible sheaves (Definition~\ref{def:native-numerical-ample-cone}) and
$\operatorname{Nef}(X)$ for the cone of classes of nonnegative degree on every closed integral
curve (Theorem~\ref{thm:native-numerical-nef-cone}).

This chapter states the objects of Fang--Han--Xia, Theorem~1.1 and Corollary~1.2, on $X$
itself: the analytification, Boucksom--Jonsson's $\theta$-psh functions for
$\theta\in N^1(X)_\R$, and the envelope property. Boucksom--Jonsson work in the finite
dimensional space $\operatorname{Num}(X)=N^1(X)_\R$, in which the ample cone is open. Finite
dimensionality of $N^1(X)_\R$ is not used here. Instead, convergence of classes is convergence
in the order of the ample cone (Definition~\ref{def:var-ample-convergence}); it is the
convergence of the vector space topology when $N^1(X)_\R$ is finite dimensional
(Theorem~\ref{thm:var-ample-convergence-fd}), and it is the convergence for which the use of
the openness of the ample cone in Boucksom--Jonsson and Fang--Han--Xia holds without further
hypothesis (Lemma~\ref{lem:var-ample-convergence-openness}).

\subsection{The analytification and convergence of numerical classes}

\begin{definition}[The analytification]
\label{def:var-analytification}
\lean{EnvelopeProperty.Variety.An}
\leanok
\uses{scheme:relative-valued-points}
\source{bj22-trivval:sec:triv}
The analytification $X^{\mathrm{an}}$ is the space of pairs $(x,|\cdot|)$ of a point $x\in X$
and an absolute value on its residue field $\kappa(x)$ that is trivial on $k$, with the
Berkovich topology (the coarsest topology for which the supports and the evaluations of
sections are continuous).
\end{definition}

\begin{definition}[Convergence in the ample-cone order]
\label{def:var-ample-convergence}
\lean{EnvelopeProperty.Variety.AmpleTendsto}
\leanok
\uses{def:native-numerical-ample-cone}
\source{bj22-trivval:defi:psh}
Let $(c_j)_{j\in J}$ be a family in $N^1(X)_\R$, $\mathcal F$ a filter on $J$ and
$\theta\in N^1(X)_\R$. We say that $c_j\to\theta$ along $\mathcal F$ if for every
$\eta\in\operatorname{Amp}(X)$, eventually along $\mathcal F$ both $\eta+(c_j-\theta)$ and
$\eta-(c_j-\theta)$ lie in $\operatorname{Amp}(X)$.
\end{definition}

\begin{lemma}[Limits]
\label{lem:var-ample-convergence-algebra}
\lean{EnvelopeProperty.Variety.ampleTendsto_const, EnvelopeProperty.Variety.AmpleTendsto.mono_left,
EnvelopeProperty.Variety.AmpleTendsto.comp, EnvelopeProperty.Variety.AmpleTendsto.congr',
EnvelopeProperty.Variety.AmpleTendsto.add, EnvelopeProperty.Variety.AmpleTendsto.neg,
EnvelopeProperty.Variety.AmpleTendsto.sub, EnvelopeProperty.Variety.AmpleTendsto.const_smul}
\leanok
\uses{def:var-ample-convergence}
Constant families converge to their value; convergence is preserved by finer filters, by
composition with maps tending to the filter, and by eventual equality. If $c_j\to\theta$ and
$c'_j\to\theta'$, then $c_j+c'_j\to\theta+\theta'$, $-c_j\to-\theta$, $c_j-c'_j\to\theta-\theta'$
and $rc_j\to r\theta$ for every $r\in\R$.
\end{lemma}
\begin{proof}
For a constant family, $\eta\pm(\theta-\theta)=\eta$. The statements on filters are immediate
from the definition. For the sum, let $\eta$ be ample; then $\eta/2$ is ample, eventually
$\eta/2\pm(c_j-\theta)$ and $\eta/2\pm(c'_j-\theta')$ are ample, and
$\eta\pm(c_j+c'_j-\theta-\theta')$ is the sum of two of these classes. For the negative, the two
conditions are exchanged. The difference is the sum with the negative. For $r>0$, apply the
definition to the ample class $r^{-1}\eta$ and multiply by $r$; for $r<0$, do the same with
$-r$ and exchange the two conditions; for $r=0$ the family is constant.
\end{proof}

\begin{lemma}[Openness of the ample cone along a convergent family]
\label{lem:var-ample-convergence-openness}
\lean{EnvelopeProperty.Variety.AmpleTendsto.eventually_sub_mem,
EnvelopeProperty.Variety.AmpleTendsto.eventually_add_sub_mem}
\leanok
\uses{def:var-ample-convergence}
\source{fhx-envelope-20261028:lem:finite-class-promotion}
If $c_j\to\theta$ and $\eta$ is ample, then eventually $\theta+\eta-c_j$ and $c_j+\eta-\theta$
are ample.
\end{lemma}
\begin{proof}
These classes are $\eta-(c_j-\theta)$ and $\eta+(c_j-\theta)$.
\end{proof}

\begin{lemma}[Order convergence implies convergence of all curve degrees]
\label{lem:var-ample-convergence-curve-weak}
\lean{EnvelopeProperty.Variety.AmpleTendsto.tendsto, EnvelopeProperty.Variety.AmpleTendsto.unique}
\leanok
\uses{def:var-ample-convergence,thm:native-numerical-curve-weak-convergence,lem:native-projectivity-geometric-ample-line,thm:native-numerical-ample-classes,thm:native-numerical-positive-ample-sums}
If $c_j\to\theta$, then $\deg_C c_j\to\deg_C\theta$ for every closed integral curve
$C\subset X$, i.e.\ $c_j\to\theta$ in the curve-weak topology of $N^1(X)_\R$. In particular
limits along a proper filter are unique.
\end{lemma}
\begin{proof}
Fix $C$ and $\varepsilon>0$. There is an ample invertible sheaf $A$
(Lemma~\ref{lem:native-projectivity-geometric-ample-line}), and $a=\deg_C A>0$
(Theorem~\ref{thm:native-numerical-ample-classes}). The class $\eta=\frac{\varepsilon}{2a}[A]$
is ample with $\deg_C\eta=\varepsilon/2$. Eventually $\eta\pm(c_j-\theta)$ are ample, and ample
classes have positive degree on $C$ (Theorem~\ref{thm:native-numerical-positive-ample-sums}).
Hence $|\deg_C c_j-\deg_C\theta|<\varepsilon/2$ eventually. Convergence of all curve degrees
is convergence in the curve-weak topology
(Theorem~\ref{thm:native-numerical-curve-weak-convergence}), which is Hausdorff.
\end{proof}

\begin{lemma}[Convergence along a finite line family]
\label{lem:var-ample-convergence-line-family}
\lean{EnvelopeProperty.Variety.ampleTendsto_lineFamilyMap}
\leanok
\uses{def:var-ample-convergence,thm:native-numerical-ample-family-openness,thm:native-numerical-line-family-maps}
Let $X$ be integral, $(L_i)_{i\in I}$ a finite family of invertible sheaves and $a_j\in\R^I$
with $a_j\to b$ in $\R^I$. Then $\sum_i a_{j,i}[L_i]\to\sum_i b_i[L_i]$.
\end{lemma}
\begin{proof}
Let $\eta$ be ample. By Theorem~\ref{thm:native-numerical-ample-family-openness} there is
$\varepsilon>0$ such that $\eta+\sum_ix_i[L_i]$ is ample whenever $|x_i|<\varepsilon$ for all
$i$. Eventually $|a_{j,i}-b_i|<\varepsilon$ for all $i$; apply this to $x=a_j-b$ and
$x=b-a_j$.
\end{proof}

\subsection{Ampleness is numerical}

\begin{theorem}[Kleiman's criterion]
\label{thm:var-numerical-ampleness}
\lean{EnvelopeProperty.Variety.isAmple_of_lineClass_mem_ampleCone,
EnvelopeProperty.Variety.isAmple_iff_lineClass_mem_ampleCone,
EnvelopeProperty.Variety.isAmple_of_lineClass_eq}
\leanok
\uses{thm:native-numerical-kleiman-from-nef-ample,def:native-numerical-ample-cone}
Let $X$ be integral and $T$ an invertible sheaf. Then $T$ is ample if and only if
$[T]\in\operatorname{Amp}(X)$. In particular ampleness depends only on the numerical class.
\end{theorem}
\begin{proof}
Ample sheaves have ample classes by definition of $\operatorname{Amp}(X)$. The converse is
Theorem~\ref{thm:native-numerical-kleiman-from-nef-ample}: write $[T]=\sum_i c_i[A_i]$ with
finitely many ample $A_i$ and $c_i>0$, choose rationals $c_i/2<q_i<c_i$ and an ample $B$ and
$d>0$ with $\sum_iq_i[A_i]=d^{-1}[B]$; then $N=T^{\otimes d}\otimes B^\vee$ has the nef class
$d\sum_i(c_i-q_i)[A_i]$, so $N\otimes B\cong T^{\otimes d}$ is ample by nef plus ample is ample,
and so is $T$. If $[T]=[T']$ and $T$ is ample, then $[T']$ is ample, hence $T'$ is ample.
\end{proof}

\subsection{Psh functions and the envelope property}

A \emph{Fubini--Study datum} $D=(M,m,(s_j,\lambda_j)_j)$ consists of an invertible sheaf $M$,
an integer $m>0$ and finitely many global sections $s_j$ of $M$ with real weights $\lambda_j$; it
represents the rational line bundle $L=M/m\in\operatorname{Pic}(X)_\Q$, of class
$[D]=m^{-1}[M]$, and has the function $\varphi_D=m^{-1}\max_j(\log|s_j|+\lambda_j)$ on
$X^{\mathrm{an}}$ (Definition~\ref{def:native-numerical-whole-section-fs-data}). It is
\emph{nonzero} if some $s_j\ne0$; on an integral $X$ this is equivalent to
$\varphi_D\not\equiv-\infty$ (Theorem~\ref{thm:native-numerical-fs-nonvanishing}). The functions
$\varphi_D$ of nonzero data with $[D]=\theta$ form the set $\mathcal H^{\mathrm{gf}}_\R(\theta)$
of generically finite Fubini--Study functions of class $\theta$
(Definition~\ref{def:native-numerical-whole-section-fs}).

\begin{definition}[$\theta$-psh functions]
\label{def:var-psh}
\lean{EnvelopeProperty.Variety.IsPSH}
\leanok
\uses{def:var-analytification,def:var-ample-convergence,def:native-numerical-whole-section-fs-data}
\source{bj22-trivval:defi:psh,bj22-trivval:defi:FS}
Let $X$ be integral and $\theta\in N^1(X)_\R$. A function
$\varphi:X^{\mathrm{an}}\to[-\infty,+\infty]$ is \emph{$\theta$-psh} if it is not identically
$-\infty$, upper semicontinuous, and the pointwise limit (infimum) of a decreasing net
$(\varphi_{D_j})_{j\in J}$, indexed by a nonempty directed preordered set, of functions of
nonzero Fubini--Study data with $[D_j]\to\theta$ in the sense of
Definition~\ref{def:var-ample-convergence}. We write $\operatorname{PSH}(X,\theta)$ for the set
of $\theta$-psh functions.
\end{definition}

\begin{remark}[Comparison with Boucksom--Jonsson]
\label{rmk:var-psh-source}
\uses{def:var-psh,thm:var-ample-convergence-fd}
Boucksom--Jonsson, Definition~4.1, requires $\lim_jc_1(L_j)=\theta$ in the finite dimensional
space $\operatorname{Num}(X)$. Definition~\ref{def:var-psh} uses the convergence of
Definition~\ref{def:var-ample-convergence}, which agrees with it whenever $N^1(X)_\R$ is finite
dimensional (Corollary~\ref{cor:var-psh-curve-weak-fd}). The index sets of the nets are taken in
a fixed universe, a size convention.
\end{remark}

\begin{definition}[The envelope property]
\label{def:var-ep}
\lean{EnvelopeProperty.Variety.HasEnvelopeProperty}
\leanok
\uses{def:var-psh}
\source{fhx-envelope-20261028:thm:intro-main,bj22-trivval:sec:envprop}
A class $\theta\in N^1(X)_\R$ has the \emph{envelope property} $\operatorname{EP}_X(\theta)$ if
for every nonempty set $S\subset\operatorname{PSH}(X,\theta)$ that is uniformly bounded above,
the upper semicontinuous regularization $(\sup S)^\star$ of its pointwise supremum is
$\theta$-psh.
\end{definition}

\begin{lemma}[Fubini--Study functions are psh]
\label{lem:var-fs-psh}
\lean{EnvelopeProperty.Variety.isPSH_of_isFS, EnvelopeProperty.Variety.isPSH_eval}
\leanok
\uses{def:var-psh,lem:var-ample-convergence-algebra,thm:native-numerical-fs-generically-finite,thm:native-numerical-fs-continuous}
If $D$ is a nonzero Fubini--Study datum, then $\varphi_D\in\operatorname{PSH}(X,[D])$.
\end{lemma}
\begin{proof}
The function $\varphi_D$ is continuous (Theorem~\ref{thm:native-numerical-fs-continuous}), hence
upper semicontinuous, and not identically $-\infty$
(Theorem~\ref{thm:native-numerical-fs-generically-finite}). The constant net $D_j=D$ over a
one-point index set is decreasing with infimum $\varphi_D$, and constant families converge.
\end{proof}

\begin{lemma}[Elementary properties]
\label{lem:var-psh-properties}
\lean{EnvelopeProperty.Variety.IsPSH.exists_ne_bot, EnvelopeProperty.Variety.IsPSH.upperSemicontinuous,
EnvelopeProperty.Variety.IsPSH.ne_top, EnvelopeProperty.Variety.IsPSH.isCurveWeakPSH}
\leanok
\uses{def:var-psh,lem:var-ample-convergence-curve-weak,thm:native-numerical-fs-ne-top,def:native-numerical-curve-weak-psh}
A $\theta$-psh function takes values in $[-\infty,+\infty)$, and it is psh for the curve-weak
topology of $N^1(X)_\R$, i.e.\ it satisfies Definition~\ref{def:var-psh} with
$[D_j]\to\theta$ replaced by convergence of all curve degrees.
\end{lemma}
\begin{proof}
Each $\varphi_{D_j}$ is finite or $-\infty$ everywhere
(Theorem~\ref{thm:native-numerical-fs-ne-top}), and $\varphi\le\varphi_{D_j}$ for one index
$j$. The second statement is Lemma~\ref{lem:var-ample-convergence-curve-weak}.
\end{proof}

\begin{lemma}[Transport of the envelope property along a homeomorphism]
\label{lem:var-ep-transport}
\lean{EnvelopeProperty.Variety.hasEnvelopeProperty_iff_of_comp_homeomorph,
EnvelopeProperty.Variety.hasEnvelopeProperty_iff_fsSystem}
\leanok
\uses{def:var-ep,def:gl-psh}
Let $e:B\to X^{\mathrm{an}}$ be a homeomorphism and $\mathcal P$ a class of functions on $B$
such that $u\circ e\in\mathcal P$ if and only if $u\in\operatorname{PSH}(X,\theta)$, for all
$u:X^{\mathrm{an}}\to[-\infty,+\infty]$. Then $\mathcal P$ has the envelope property (every
nonempty subset bounded above has its regularized supremum in $\mathcal P$) if and only if
$\theta$ has the envelope property. In particular this applies to the $\theta_F$-psh functions
of a Fubini--Study system $F$ on $B$, such as the system of a model
(Definition~\ref{def:gl-psh}).
\end{lemma}
\begin{proof}
Composition with $e$ is a bijection between functions on $X^{\mathrm{an}}$ and on $B$ that
preserves pointwise suprema, upper bounds and, since $e$ is a homeomorphism, upper
semicontinuous regularizations: $(f\circ e)^\star=f^\star\circ e$. Given a nonempty set
$S\subset\operatorname{PSH}(X,\theta)$ bounded above, the set $\{u\circ e:u\in S\}$ lies in
$\mathcal P$, is nonempty and bounded above, and its regularized supremum is
$(\sup S)^\star\circ e$; so the envelope property of $\mathcal P$ gives
$(\sup S)^\star\in\operatorname{PSH}(X,\theta)$. The converse is the same argument with
$e^{-1}$.
\end{proof}

\begin{lemma}[The envelope property for indexed families]
\label{lem:var-ep-family}
\lean{EnvelopeProperty.Variety.HasEnvelopeProperty.isPSH_uscReg_iSup}
\leanok
\uses{def:var-ep}
If $\theta$ has the envelope property and $(\varphi_i)_{i\in I}$ is a nonempty family in
$\operatorname{PSH}(X,\theta)$ uniformly bounded above, then $(\sup_i\varphi_i)^\star$ is
$\theta$-psh.
\end{lemma}
\begin{proof}
Apply the definition to the set $\{\varphi_i\}$, whose pointwise supremum is
$\sup_i\varphi_i$.
\end{proof}

\subsection{Promotion of Fubini--Study functions}

\begin{lemma}[Finite class promotion]
\label{lem:var-class-promotion}
\lean{EnvelopeProperty.NativeNumericalFubiniStudy.FSData.exists_promotion,
EnvelopeProperty.NativeNumericalFubiniStudy.FSData.isFS_of_ample_sub}
\leanok
\uses{thm:native-numerical-kleiman-from-nef-ample,lem:native-numerical-ray-from-ample-class,thm:native-numerical-fs-from-ray,thm:native-numerical-generated-line-sums,lem:native-numerical-line-tensor-power}
\source{fhx-envelope-20261028:lem:finite-class-promotion}
Let $X$ be integral, $D=(L,m,(s_j,\lambda_j))$ a nonzero Fubini--Study datum, $M$ an invertible
sheaf and $n>0$ such that $n^{-1}[M]-[D]$ is ample. Then there are $t>0$ and a nonzero datum $E$
with line $M^{\otimes t}$, denominator $nt$ and $\varphi_E=\varphi_D$. In particular
$\varphi_D\in\mathcal H^{\mathrm{gf}}_\R(n^{-1}[M])$, the class of the rational line bundle $M/n$.
\end{lemma}
\begin{proof}
By Kleiman's criterion (Theorem~\ref{thm:native-numerical-kleiman-from-nef-ample}) and
Lemma~\ref{lem:native-numerical-ray-from-ample-class} with $b=1/n$, the function $\varphi_D$
lies in the fixed-ray class $\mathcal H(M,1/n)$: there is a datum with denominator $m'>0$, degree
$t$ with $t=m'/n$, nonzero sections $s'_j$ of $M^{\otimes t}$ and weights $\lambda'_j$, whose
function is $\varphi_D$. Concretely, the invertible sheaf
$P=M^{\otimes m}\otimes(L^{\otimes n})^\vee$ has class $mn(n^{-1}[M]-[D])$, so it is ample;
a power $P^{\otimes s}$ has finitely many global generators $(t_c)$, all of norm $\le1$ and one
of norm $1$ at every point of $X^{\mathrm{an}}$; and the datum is the twist of the $ns$-th power
of $D$ by the $t_c$, whose line $L^{\otimes ns}\otimes P^{\otimes s}$ is isomorphic to
$M^{\otimes ms}$. Regarding this fixed-ray datum as a varying-line datum
(Theorem~\ref{thm:native-numerical-fs-from-ray}) gives $E$ with line $M^{\otimes t}$ and
denominator $m'=nt$; $t>0$ because $m'>0$. The class of $E$ is
$(nt)^{-1}t[M]=n^{-1}[M]$ (Theorem~\ref{thm:native-numerical-generated-line-sums},
Lemma~\ref{lem:native-numerical-line-tensor-power}).
\end{proof}

\begin{remark}[Relation to the source]
\label{rmk:var-class-promotion-source}
\uses{lem:var-class-promotion}
Fang--Han--Xia state the promotion for a finite surjective morphism $h:T\to X$ and a sequence
of classes; Lemma~\ref{lem:var-class-promotion} is the step of its proof for $h$ the identity
and one function, with the rational line bundle $M/n$ prescribed. The choice of the classes
$n^{-1}[M]$ and the use of the openness of the ample cone
(Lemma~\ref{lem:var-ample-convergence-openness}) are left to the application.
\end{remark}

\subsection{Finite dimensional numerical spaces}

In this subsection $X$ is integral and $N^1(X)_\R$ is assumed to be finite dimensional; this is
a hypothesis of the statements below, which hold unconditionally once the finite dimensionality
of $N^1(X)_\R$ is known.

\begin{lemma}[The ample cone is open]
\label{lem:var-ample-cone-open-fd}
\lean{EnvelopeProperty.Variety.isOpen_ampleCone}
\leanok
\uses{thm:native-numerical-line-family-maps,thm:native-numerical-ample-family-openness,thm:native-numerical-curve-weak-vector-topology}
If $N^1(X)_\R$ is finite dimensional, then $\operatorname{Amp}(X)$ is open in $N^1(X)_\R$ with
its curve-weak topology.
\end{lemma}
\begin{proof}
Choose a basis $(e_i)$ of $N^1(X)_\R$ and write each $e_i$ as a finite real combination of
classes of invertible sheaves (Theorem~\ref{thm:native-numerical-line-family-maps}). Collect
all these sheaves into one finite family $(L_s)_{s\in S}$. Its coefficient map
$F:\R^S\to N^1(X)_\R$ contains every $[L_s]$, hence every $e_i$, in its image, so $F$ is
surjective. The preimage $F^{-1}(\operatorname{Amp}(X))$ is open
(Theorem~\ref{thm:native-numerical-ample-family-openness}), and a surjective linear map onto a
finite dimensional Hausdorff topological vector space is open
(Theorem~\ref{thm:native-numerical-curve-weak-vector-topology}). Hence
$\operatorname{Amp}(X)=F(F^{-1}(\operatorname{Amp}(X)))$ is open.
\end{proof}

\begin{theorem}[Order convergence is topological convergence]
\label{thm:var-ample-convergence-fd}
\lean{EnvelopeProperty.Variety.ampleTendsto_iff_tendsto}
\leanok
\uses{lem:var-ample-cone-open-fd,lem:var-ample-convergence-curve-weak}
If $N^1(X)_\R$ is finite dimensional, then $c_j\to\theta$ in the sense of
Definition~\ref{def:var-ample-convergence} if and only if $c_j\to\theta$ in the topology of
$N^1(X)_\R$.
\end{theorem}
\begin{proof}
One direction is Lemma~\ref{lem:var-ample-convergence-curve-weak}. Conversely, if
$c_j\to\theta$, then $c_j-\theta\to0$. For ample $\eta$, the set
$\{x:\eta+x\in\operatorname{Amp}(X),\ \eta-x\in\operatorname{Amp}(X)\}$ is open
(Lemma~\ref{lem:var-ample-cone-open-fd}) and contains $0$, so it eventually contains
$c_j-\theta$.
\end{proof}

\begin{corollary}[Comparison with Boucksom--Jonsson's definition]
\label{cor:var-psh-curve-weak-fd}
\lean{EnvelopeProperty.Variety.isPSH_iff_isCurveWeakPSH}
\leanok
\uses{thm:var-ample-convergence-fd,def:var-psh,def:native-numerical-curve-weak-psh}
\source{bj22-trivval:defi:psh}
If $N^1(X)_\R$ is finite dimensional, then $\operatorname{PSH}(X,\theta)$ is the set of
functions of Boucksom--Jonsson, Definition~4.1, with convergence of the classes in the topology
of $N^1(X)_\R$.
\end{corollary}
\begin{proof}
The two definitions differ only in the convergence of the classes $[D_j]$, and the two
convergences agree by Theorem~\ref{thm:var-ample-convergence-fd}.
\end{proof}
